% References -- all sources cited in Chapters 1--23, consolidated and
% listed alphabetically by first author. Generated from the chapters'
% former per-chapter reference lists.
\chapter*{References}
\addcontentsline{toc}{chapter}{References}
\markboth{References}{References}
\label{references}

\begingroup
\small
\sloppy
\setlength{\parindent}{0pt}
\setlength{\parskip}{3pt}

\hangindent=1.5em\hangafter=1 Abdi, F. and Ranaldo, A. (2017). A Simple Estimation of Bid-Ask Spreads from Daily Close, High, and Low Prices. \emph{Review of Financial Studies} 30(12): 4437--4480.

\hangindent=1.5em\hangafter=1 Acerbi, C. and Sz\'ekely, B. (2014). Back-Testing Expected Shortfall. \emph{Risk}, December 2014.

\hangindent=1.5em\hangafter=1 Adrian, T., Crump, R. K. and Moench, E. (2013). Pricing the Term Structure with Linear Regressions. \emph{Journal of Financial Economics} 110(1): 110--138.

\hangindent=1.5em\hangafter=1 Almgren, R., Thum, C., Hauptmann, E. and Li, H. (2005). Direct Estimation of Equity Market Impact. \emph{Risk} 18(7): 58--62.

\hangindent=1.5em\hangafter=1 American Finance Association (2021). Answers to FAQ about the Recent Retraction of an Article in the JF. July 10.

\hangindent=1.5em\hangafter=1 Amihud, Y. (2002). Illiquidity and Stock Returns: Cross-Section and Time-Series Effects. \emph{Journal of Financial Markets} 5(1): 31--56.

\hangindent=1.5em\hangafter=1 Amihud, Y. and Hurvich, C. M. (2004). Predictive Regressions: A Reduced-Bias Estimation Method. \emph{Journal of Financial and Quantitative Analysis} 39(4): 813--841.

\hangindent=1.5em\hangafter=1 Andersen, T. G. and Bollerslev, T. (1998). Answering the Skeptics: Yes, Standard Volatility Models Do Provide Accurate Forecasts. \emph{International Economic Review} 39(4): 885--905.

\hangindent=1.5em\hangafter=1 Andersen, T. G., Bollerslev, T., Diebold, F. X. and Labys, P. (2003). Modeling and Forecasting Realized Volatility. \emph{Econometrica} 71(2): 579--625.

\hangindent=1.5em\hangafter=1 Ang, A. and Bekaert, G. (2007). Stock Return Predictability: Is It There? \emph{Review of Financial Studies} 20(3): 651--707.

\hangindent=1.5em\hangafter=1 Ang, A. and Chen, J. (2002). Asymmetric Correlations of Equity Portfolios. \emph{Journal of Financial Economics} 63(3): 443--494.

\hangindent=1.5em\hangafter=1 Ang, A., Hodrick, R. J., Xing, Y. and Zhang, X. (2006). The Cross-Section of Volatility and Expected Returns. \emph{Journal of Finance} 61(1): 259--299.

\hangindent=1.5em\hangafter=1 Ang, A., Hodrick, R. J., Xing, Y. and Zhang, X. (2009). High Idiosyncratic Volatility and Low Returns: International and Further U.S. Evidence. \emph{Journal of Financial Economics} 91(1): 1--23.

\hangindent=1.5em\hangafter=1 Angrist, J. D. and Pischke, J.-S. (2009). \emph{Mostly Harmless Econometrics: An Empiricist's Companion}. Princeton: Princeton University Press.

\hangindent=1.5em\hangafter=1 Ansari, A. F., et al.\ (2024). Chronos: Learning the Language of Time Series. \emph{Transactions on Machine Learning Research}.

\hangindent=1.5em\hangafter=1 Arnott, R. D., Harvey, C. R., Kalesnik, V. and Linnainmaa, J. T. (2021). Reports of Value's Death May Be Greatly Exaggerated. \emph{Financial Analysts Journal} 77(1): 44--67.

\hangindent=1.5em\hangafter=1 Artzner, P., Delbaen, F., Eber, J.-M. and Heath, D. (1999). Coherent Measures of Risk. \emph{Mathematical Finance} 9(3): 203--228.

\hangindent=1.5em\hangafter=1 Asness, C., Frazzini, A., Gormsen, N. J. and Pedersen, L. H. (2020). Betting against Correlation: Testing Theories of the Low-Risk Effect. \emph{Journal of Financial Economics} 135(3): 629--652.

\hangindent=1.5em\hangafter=1 Asness, C. S. (2003). Fight the Fed Model. \emph{Journal of Portfolio Management} 30(1): 11--24.

\hangindent=1.5em\hangafter=1 Asness, C. S., Moskowitz, T. J. and Pedersen, L. H. (2013). Value and Momentum Everywhere. \emph{Journal of Finance} 68(3): 929--985.

\hangindent=1.5em\hangafter=1 Avramov, D., Cheng, S. and Metzker, L. (2023). Machine Learning vs.\ Economic Restrictions: Evidence from Stock Return Predictability. \emph{Management Science} 69(5): 2587--2619.

\hangindent=1.5em\hangafter=1 Bailey, D. H. and L\'opez de Prado, M. (2012). The Sharpe Ratio Efficient Frontier. \emph{Journal of Risk} 15(2): 3--44.

\hangindent=1.5em\hangafter=1 Bailey, D. H. and L\'opez de Prado, M. (2014). The Deflated Sharpe Ratio: Correcting for Selection Bias, Backtest Overfitting, and Non-Normality. \emph{Journal of Portfolio Management} 40(5): 94--107.

\hangindent=1.5em\hangafter=1 Bailey, D. H., Borwein, J. M., L\'opez de Prado, M. and Zhu, Q. J. (2014). Pseudo-Mathematics and Financial Charlatanism: The Effects of Backtest Overfitting on Out-of-Sample Performance. \emph{Notices of the American Mathematical Society} 61(5): 458--471.

\hangindent=1.5em\hangafter=1 Bailey, D. H., Borwein, J. M., L\'opez de Prado, M. and Zhu, Q. J. (2017). The Probability of Backtest Overfitting. \emph{Journal of Computational Finance} 20(4): 39--69.

\hangindent=1.5em\hangafter=1 Bajgrowicz, P. and Scaillet, O. (2012). Technical Trading Revisited: False Discoveries, Persistence Tests, and Transaction Costs. \emph{Journal of Financial Economics} 106(3): 473--491.

\hangindent=1.5em\hangafter=1 Baker, M., Bradley, B. and Wurgler, J. (2011). Benchmarks as Limits to Arbitrage: Understanding the Low-Volatility Anomaly. \emph{Financial Analysts Journal} 67(1): 40--54.

\hangindent=1.5em\hangafter=1 Bakshi, G. and Kapadia, N. (2003). Delta-Hedged Gains and the Negative Market Volatility Risk Premium. \emph{Review of Financial Studies} 16(2): 527--566.

\hangindent=1.5em\hangafter=1 Bakshi, G., Cao, C. and Chen, Z. (1997). Empirical Performance of Alternative Option Pricing Models. \emph{Journal of Finance} 52(5): 2003--2049.

\hangindent=1.5em\hangafter=1 Bali, T. G., Cakici, N. and Whitelaw, R. F. (2011). Maxing Out: Stocks as Lotteries and the Cross-Section of Expected Returns. \emph{Journal of Financial Economics} 99(2): 427--446.

\hangindent=1.5em\hangafter=1 Banz, R. W. (1981). The Relationship Between Return and Market Value of Common Stocks. \emph{Journal of Financial Economics} 9(1): 3--18.

\hangindent=1.5em\hangafter=1 Barber, B. M. and Odean, T. (2000). Trading Is Hazardous to Your Wealth: The Common Stock Investment Performance of Individual Investors. \emph{Journal of Finance} 55(2): 773--806.

\hangindent=1.5em\hangafter=1 Barber, B. M., Huang, X., Odean, T. and Schwarz, C. (2022). Attention-Induced Trading and Returns: Evidence from Robinhood Users. \emph{Journal of Finance} 77(6): 3141--3190.

\hangindent=1.5em\hangafter=1 Barber, B. M., Lee, Y.-T., Liu, Y.-J. and Odean, T. (2009). Just How Much Do Individual Investors Lose by Trading? \emph{Review of Financial Studies} 22(2): 609--632.

\hangindent=1.5em\hangafter=1 Barberis, N. and Huang, M. (2008). Stocks as Lotteries: The Implications of Probability Weighting for Security Prices. \emph{American Economic Review} 98(5): 2066--2100.

\hangindent=1.5em\hangafter=1 Barberis, N. and Thaler, R. (2003). A Survey of Behavioral Finance. In \emph{Handbook of the Economics of Finance}, vol.\ 1B, 1053--1128.

\hangindent=1.5em\hangafter=1 Barberis, N., Shleifer, A. and Vishny, R. (1998). A Model of Investor Sentiment. \emph{Journal of Financial Economics} 49(3): 307--343.

\hangindent=1.5em\hangafter=1 Barone-Adesi, G., Giannopoulos, K. and Vosper, L. (1999). VaR without Correlations for Portfolios of Derivative Securities. \emph{Journal of Futures Markets} 19(5): 583--602.

\hangindent=1.5em\hangafter=1 Barroso, P. and Santa-Clara, P. (2015). Momentum Has Its Moments. \emph{Journal of Financial Economics} 116(1): 111--120.

\hangindent=1.5em\hangafter=1 Basel Committee on Banking Supervision (1996). \emph{Supervisory Framework for the Use of ``Backtesting'' in Conjunction with the Internal Models Approach to Market Risk Capital Requirements}. Basel: Bank for International Settlements.

\hangindent=1.5em\hangafter=1 Basel Committee on Banking Supervision (2009). \emph{Revisions to the Basel II Market Risk Framework}. Basel: Bank for International Settlements.

\hangindent=1.5em\hangafter=1 Basel Committee on Banking Supervision (2016). \emph{Minimum Capital Requirements for Market Risk}. Basel: Bank for International Settlements.

\hangindent=1.5em\hangafter=1 Basu, S. (1977). Investment Performance of Common Stocks in Relation to Their Price-Earnings Ratios: A Test of the Efficient Market Hypothesis. \emph{Journal of Finance} 32(3): 663--682.

\hangindent=1.5em\hangafter=1 Bauer, M. D. and Hamilton, J. D. (2018). Robust Bond Risk Premia. \emph{Review of Financial Studies} 31(2): 399--448.

\hangindent=1.5em\hangafter=1 Belkin, M., Hsu, D., Ma, S. and Mandal, S. (2019). Reconciling Modern Machine-Learning Practice and the Classical Bias--Variance Trade-Off. \emph{Proceedings of the National Academy of Sciences} 116(32): 15849--15854.

\hangindent=1.5em\hangafter=1 Bengen, W. P. (1994). Determining Withdrawal Rates Using Historical Data. \emph{Journal of Financial Planning} 7(4): 171--180.

\hangindent=1.5em\hangafter=1 Benjamin, D. J., Berger, J. O., Johannesson, M., Nosek, B. A., Wagenmakers, E.-J., et al. (2018). Redefine Statistical Significance. \emph{Nature Human Behaviour} 2(1): 6--10.

\hangindent=1.5em\hangafter=1 Benjamini, Y. and Hochberg, Y. (1995). Controlling the False Discovery Rate: A Practical and Powerful Approach to Multiple Testing. \emph{Journal of the Royal Statistical Society: Series B} 57(1): 289--300.

\hangindent=1.5em\hangafter=1 Benjamini, Y. and Yekutieli, D. (2001). The Control of the False Discovery Rate in Multiple Testing under Dependency. \emph{Annals of Statistics} 29(4): 1165--1188.

\hangindent=1.5em\hangafter=1 Berk, J. B. and Green, R. C. (2004). Mutual Fund Flows and Performance in Rational Markets. \emph{Journal of Political Economy} 112(6): 1269--1295.

\hangindent=1.5em\hangafter=1 Berkowitz, J. and O'Brien, J. (2002). How Accurate Are Value-at-Risk Models at Commercial Banks? \emph{Journal of Finance} 57(3): 1093--1111.

\hangindent=1.5em\hangafter=1 Bernard, V. L. and Thomas, J. K. (1989). Post-Earnings-Announcement Drift: Delayed Price Response or Risk Premium? \emph{Journal of Accounting Research} 27(Supplement): 1--36.

\hangindent=1.5em\hangafter=1 Bianchi, D., B\"uchner, M. and Tamoni, A. (2021). Bond Risk Premiums with Machine Learning. \emph{Review of Financial Studies} 34(2): 1046--1089.

\hangindent=1.5em\hangafter=1 Black, F. (1972). Capital Market Equilibrium with Restricted Borrowing. \emph{Journal of Business} 45(3): 444--455.

\hangindent=1.5em\hangafter=1 Black, F. (1976). Studies of Stock Price Volatility Changes. \emph{Proceedings of the 1976 Meetings of the American Statistical Association, Business and Economic Statistics Section}: 177--181.

\hangindent=1.5em\hangafter=1 Black, F. and Litterman, R. (1992). Global Portfolio Optimization. \emph{Financial Analysts Journal} 48(5): 28--43.

\hangindent=1.5em\hangafter=1 Black, F. and Scholes, M. (1972). The Valuation of Option Contracts and a Test of Market Efficiency. \emph{Journal of Finance} 27(2): 399--417.

\hangindent=1.5em\hangafter=1 Black, F. and Scholes, M. (1973). The Pricing of Options and Corporate Liabilities. \emph{Journal of Political Economy} 81(3): 637--654.

\hangindent=1.5em\hangafter=1 Black, F., Jensen, M. C. and Scholes, M. (1972). The Capital Asset Pricing Model: Some Empirical Tests. In M. C. Jensen (ed.), \emph{Studies in the Theory of Capital Markets}, 79--121. New York: Praeger.

\hangindent=1.5em\hangafter=1 Blanchett, D. (2017). Testing DALBAR's Claims about Mutual Fund Investors. \emph{Advisor Perspectives}, May 8.

\hangindent=1.5em\hangafter=1 Board of Governors of the Federal Reserve System (2021). Banks' Backtesting Exceptions during the COVID-19 Crash: Causes and Consequences. \emph{FEDS Notes}, July 8.

\hangindent=1.5em\hangafter=1 Bollen, J., Mao, H. and Zeng, X. (2011). Twitter Mood Predicts the Stock Market. \emph{Journal of Computational Science} 2(1): 1--8.

\hangindent=1.5em\hangafter=1 Bollen, N. P. B. and Whaley, R. E. (2004). Does Net Buying Pressure Affect the Shape of Implied Volatility Functions? \emph{Journal of Finance} 59(2): 711--753.

\hangindent=1.5em\hangafter=1 Bollerslev, T. (1986). Generalized Autoregressive Conditional Heteroskedasticity. \emph{Journal of Econometrics} 31(3): 307--327.

\hangindent=1.5em\hangafter=1 Bollerslev, T., Tauchen, G. and Zhou, H. (2009). Expected Stock Returns and Variance Risk Premia. \emph{Review of Financial Studies} 22(11): 4463--4492.

\hangindent=1.5em\hangafter=1 Boudoukh, J., Richardson, M. and Whitelaw, R. F. (2008). The Myth of Long-Horizon Predictability. \emph{Review of Financial Studies} 21(4): 1577--1605.

\hangindent=1.5em\hangafter=1 Boyer, B. H., Gibson, M. S. and Loretan, M. (1999). Pitfalls in Tests for Changes in Correlations. International Finance Discussion Paper 597, Board of Governors of the Federal Reserve System.

\hangindent=1.5em\hangafter=1 Breiman, L. (1996). Bagging Predictors. \emph{Machine Learning} 24(2): 123--140.

\hangindent=1.5em\hangafter=1 Breiman, L. (2001). Random Forests. \emph{Machine Learning} 45(1): 5--32.

\hangindent=1.5em\hangafter=1 Brini, A. (2026). Forecasting Realized Volatility with Time Series Foundation Models: A Comparison with Econometric Benchmarks. arXiv:2607.05291.

\hangindent=1.5em\hangafter=1 Broadie, M., Chernov, M. and Johannes, M. (2009). Understanding Index Option Returns. \emph{Review of Financial Studies} 22(11): 4493--4529.

\hangindent=1.5em\hangafter=1 Brock, W., Lakonishok, J. and LeBaron, B. (1992). Simple Technical Trading Rules and the Stochastic Properties of Stock Returns. \emph{Journal of Finance} 47(5): 1731--1764.

\hangindent=1.5em\hangafter=1 Brown, S. J., Goetzmann, W., Ibbotson, R. G. and Ross, S. A. (1992). Survivorship Bias in Performance Studies. \emph{Review of Financial Studies} 5(4): 553--580.

\hangindent=1.5em\hangafter=1 Brunnermeier, M. K., Nagel, S. and Pedersen, L. H. (2008). Carry Trades and Currency Crashes. \emph{NBER Macroeconomics Annual} 23: 313--347.

\hangindent=1.5em\hangafter=1 Burnside, C., Eichenbaum, M., Kleshchelski, I. and Rebelo, S. (2011). Do Peso Problems Explain the Returns to the Carry Trade? \emph{Review of Financial Studies} 24(3): 853--891.

\hangindent=1.5em\hangafter=1 Callan Associates (2006). \emph{An Historical Evaluation of the CBOE S\&P 500 BuyWrite Index Strategy}. Chicago: Cboe.

\hangindent=1.5em\hangafter=1 Campbell, J. Y. (1987). Stock Returns and the Term Structure. \emph{Journal of Financial Economics} 18(2): 373--399.

\hangindent=1.5em\hangafter=1 Campbell, J. Y. and Shiller, R. J. (1988). The Dividend-Price Ratio and Expectations of Future Dividends and Discount Factors. \emph{Review of Financial Studies} 1(3): 195--228.

\hangindent=1.5em\hangafter=1 Campbell, J. Y. and Shiller, R. J. (1991). Yield Spreads and Interest Rate Movements: A Bird's Eye View. \emph{Review of Economic Studies} 58(3): 495--514.

\hangindent=1.5em\hangafter=1 Campbell, J. Y. and Shiller, R. J. (1998). Valuation Ratios and the Long-Run Stock Market Outlook. \emph{Journal of Portfolio Management} 24(2): 11--26.

\hangindent=1.5em\hangafter=1 Campbell, J. Y. and Thompson, S. B. (2008). Predicting Excess Stock Returns Out of Sample: Can Anything Beat the Historical Average? \emph{Review of Financial Studies} 21(4): 1509--1531.

\hangindent=1.5em\hangafter=1 Campbell, J. Y. and Yogo, M. (2006). Efficient Tests of Stock Return Predictability. \emph{Journal of Financial Economics} 81(1): 27--60.

\hangindent=1.5em\hangafter=1 Campbell, J. Y., Sunderam, A. and Viceira, L. M. (2017). Inflation Bets or Deflation Hedges? The Changing Risks of Nominal Bonds. \emph{Critical Finance Review} 6(2): 263--301.

\hangindent=1.5em\hangafter=1 Carhart, M. M. (1997). On Persistence in Mutual Fund Performance. \emph{Journal of Finance} 52(1): 57--82.

\hangindent=1.5em\hangafter=1 Carr, P. and Wu, L. (2009). Variance Risk Premiums. \emph{Review of Financial Studies} 22(3): 1311--1341.

\hangindent=1.5em\hangafter=1 Cboe Global Markets (2025). SPX 0DTE Options Jump to 61\% Share on Retail Resurgence. \emph{Cboe Insights}, June 2.

\hangindent=1.5em\hangafter=1 Cederburg, S. and O'Doherty, M. S. (2016). Does It Pay to Bet against Beta? On the Conditional Performance of the Beta Anomaly. \emph{Journal of Finance} 71(2): 737--774.

\hangindent=1.5em\hangafter=1 Cederburg, S., O'Doherty, M. S., Wang, F. and Yan, X. (2020). On the Performance of Volatility-Managed Portfolios. \emph{Journal of Financial Economics} 138(1): 95--117.

\hangindent=1.5em\hangafter=1 Chan, L. K. C., Lakonishok, J. and Sougiannis, T. (2001). The Stock Market Valuation of Research and Development Expenditures. \emph{Journal of Finance} 56(6): 2431--2456.

\hangindent=1.5em\hangafter=1 Chen, A. Y. (2021). The Limits of p-Hacking: Some Thought Experiments. \emph{Journal of Finance} 76(5): 2447--2480.

\hangindent=1.5em\hangafter=1 Chen, A. Y. and Velikov, M. (2023). Zeroing In on the Expected Returns of Anomalies. \emph{Journal of Financial and Quantitative Analysis} 58(3): 968--1004.

\hangindent=1.5em\hangafter=1 Chen, A. Y. and Zimmermann, T. (2020). Publication Bias and the Cross-Section of Stock Returns. \emph{Review of Asset Pricing Studies} 10(2): 249--289.

\hangindent=1.5em\hangafter=1 Chen, A. Y. and Zimmermann, T. (2022). Open Source Cross-Sectional Asset Pricing. \emph{Critical Finance Review} 11(2): 207--264.

\hangindent=1.5em\hangafter=1 Chen, L., Pelger, M. and Zhu, J. (2024). Deep Learning in Asset Pricing. \emph{Management Science} 70(2): 714--750.

\hangindent=1.5em\hangafter=1 Chen, Y., Kelly, B. T. and Xiu, D. (2022). Expected Returns and Large Language Models. Working paper, University of Chicago.

\hangindent=1.5em\hangafter=1 Chopra, V. K. and Ziemba, W. T. (1993). The Effect of Errors in Means, Variances, and Covariances on Optimal Portfolio Choice. \emph{Journal of Portfolio Management} 19(2): 6--11.

\hangindent=1.5em\hangafter=1 Christoffersen, P. F. (1998). Evaluating Interval Forecasts. \emph{International Economic Review} 39(4): 841--862.

\hangindent=1.5em\hangafter=1 Chui, A. C. W., Titman, S. and Wei, K. C. J. (2010). Individualism and Momentum around the World. \emph{Journal of Finance} 65(1): 361--392.

\hangindent=1.5em\hangafter=1 Clark, T. E. and West, K. D. (2007). Approximately Normal Tests for Equal Predictive Accuracy in Nested Models. \emph{Journal of Econometrics} 138(1): 291--311.

\hangindent=1.5em\hangafter=1 Cochrane, J. H. (2008). The Dog That Did Not Bark: A Defense of Return Predictability. \emph{Review of Financial Studies} 21(4): 1533--1575.

\hangindent=1.5em\hangafter=1 Cochrane, J. H. (2011). Presidential Address: Discount Rates. \emph{Journal of Finance} 66(4): 1047--1108.

\hangindent=1.5em\hangafter=1 Cochrane, J. H. and Piazzesi, M. (2005). Bond Risk Premia. \emph{American Economic Review} 95(1): 138--160.

\hangindent=1.5em\hangafter=1 Collin-Dufresne, P., Goldstein, R. S. and Martin, J. S. (2001). The Determinants of Credit Spread Changes. \emph{Journal of Finance} 56(6): 2177--2207.

\hangindent=1.5em\hangafter=1 Cont, R. (2001). Empirical Properties of Asset Returns: Stylized Facts and Statistical Issues. \emph{Quantitative Finance} 1(2): 223--236.

\hangindent=1.5em\hangafter=1 Cooper, M. J., Gulen, H. and Schill, M. J. (2008). Asset Growth and the Cross-Section of Stock Returns. \emph{Journal of Finance} 63(4): 1609--1651.

\hangindent=1.5em\hangafter=1 Corsi, F. (2009). A Simple Approximate Long-Memory Model of Realized Volatility. \emph{Journal of Financial Econometrics} 7(2): 174--196.

\hangindent=1.5em\hangafter=1 Corwin, S. A. and Schultz, P. (2012). A Simple Way to Estimate Bid-Ask Spreads from Daily High and Low Prices. \emph{Journal of Finance} 67(2): 719--760.

\hangindent=1.5em\hangafter=1 Coval, J. D. and Shumway, T. (2001). Expected Option Returns. \emph{Journal of Finance} 56(3): 983--1009.

\hangindent=1.5em\hangafter=1 Coval, J., Jurek, J. and Stafford, E. (2009). The Economics of Structured Finance. \emph{Journal of Economic Perspectives} 23(1): 3--25.

\hangindent=1.5em\hangafter=1 DALBAR, Inc.\ (2016). \emph{Quantitative Analysis of Investor Behavior}. Boston: DALBAR.

\hangindent=1.5em\hangafter=1 Daniel, K. and Moskowitz, T. J. (2016). Momentum Crashes. \emph{Journal of Financial Economics} 122(2): 221--247.

\hangindent=1.5em\hangafter=1 Daniel, K., Hirshleifer, D. and Subrahmanyam, A. (1998). Investor Psychology and Security Market Under- and Overreactions. \emph{Journal of Finance} 53(6): 1839--1885.

\hangindent=1.5em\hangafter=1 Das, A., Kong, W., Sen, R. and Zhou, Y. (2024). A Decoder-Only Foundation Model for Time-Series Forecasting. \emph{Proceedings of the 41st International Conference on Machine Learning}.

\hangindent=1.5em\hangafter=1 D'Avolio, G. (2002). The Market for Borrowing Stock. \emph{Journal of Financial Economics} 66(2--3): 271--306.

\hangindent=1.5em\hangafter=1 De Bondt, W. F. M. and Thaler, R. (1985). Does the Stock Market Overreact? \emph{Journal of Finance} 40(3): 793--805.

\hangindent=1.5em\hangafter=1 DeMiguel, V., Garlappi, L. and Uppal, R. (2009). Optimal versus Naive Diversification: How Inefficient Is the 1/N Portfolio Strategy? \emph{Review of Financial Studies} 22(5): 1915--1953.

\hangindent=1.5em\hangafter=1 Dichev, I. D. (2007). What Are Stock Investors' Actual Historical Returns? Evidence from Dollar-Weighted Returns. \emph{American Economic Review} 97(1): 386--401.

\hangindent=1.5em\hangafter=1 Diebold, F. X. and Mariano, R. S. (1995). Comparing Predictive Accuracy. \emph{Journal of Business \& Economic Statistics} 13(3): 253--263.

\hangindent=1.5em\hangafter=1 Dimensional Fund Advisors (2023). Market Review 2022: After a Down Year, Looking to the Past as a Guide.

\hangindent=1.5em\hangafter=1 Dimson, E. (1979). Risk Measurement When Shares Are Subject to Infrequent Trading. \emph{Journal of Financial Economics} 7(2): 197--226.

\hangindent=1.5em\hangafter=1 Driessen, J., Kuiper, I., Nazliben, K. and Beilo, R. (2019). Does Interest Rate Exposure Explain the Low-Volatility Anomaly? \emph{Journal of Banking \& Finance} 103: 51--61.

\hangindent=1.5em\hangafter=1 Du, W., Tepper, A. and Verdelhan, A. (2018). Deviations from Covered Interest Rate Parity. \emph{Journal of Finance} 73(3): 915--957.

\hangindent=1.5em\hangafter=1 Efron, B. (1979). Bootstrap Methods: Another Look at the Jackknife. \emph{Annals of Statistics} 7(1): 1--26.

\hangindent=1.5em\hangafter=1 Efron, B. and Morris, C. (1975). Data Analysis Using Stein's Estimator and Its Generalizations. \emph{Journal of the American Statistical Association} 70(350): 311--319.

\hangindent=1.5em\hangafter=1 Ehsani, S. and Linnainmaa, J. T. (2022). Factor Momentum and the Momentum Factor. \emph{Journal of Finance} 77(3): 1877--1919.

\hangindent=1.5em\hangafter=1 Eisfeldt, A. L., Kim, E. T. and Papanikolaou, D. (2022). Intangible Value. \emph{Critical Finance Review} 11(2): 299--332.

\hangindent=1.5em\hangafter=1 Elton, E. J., Gruber, M. J., Agrawal, D. and Mann, C. (2001). Explaining the Rate Spread on Corporate Bonds. \emph{Journal of Finance} 56(1): 247--277.

\hangindent=1.5em\hangafter=1 Elton, E. J., Gruber, M. J. and Blake, C. R. (1996). Survivorship Bias and Mutual Fund Performance. \emph{Review of Financial Studies} 9(4): 1097--1120.

\hangindent=1.5em\hangafter=1 Engle, R. F. (1982). Autoregressive Conditional Heteroscedasticity with Estimates of the Variance of United Kingdom Inflation. \emph{Econometrica} 50(4): 987--1007.

\hangindent=1.5em\hangafter=1 Faber, M. T. (2007). A Quantitative Approach to Tactical Asset Allocation. \emph{Journal of Wealth Management} 9(4): 69--79.

\hangindent=1.5em\hangafter=1 Fama, E. F. (1965). The Behavior of Stock-Market Prices. \emph{Journal of Business} 38(1): 34--105.

\hangindent=1.5em\hangafter=1 Fama, E. F. (1970). Efficient Capital Markets: A Review of Theory and Empirical Work. \emph{Journal of Finance} 25(2): 383--417.

\hangindent=1.5em\hangafter=1 Fama, E. F. (1984). Forward and Spot Exchange Rates. \emph{Journal of Monetary Economics} 14(3): 319--338.

\hangindent=1.5em\hangafter=1 Fama, E. F. and Bliss, R. R. (1987). The Information in Long-Maturity Forward Rates. \emph{American Economic Review} 77(4): 680--692.

\hangindent=1.5em\hangafter=1 Fama, E. F. and French, K. R. (1988). Dividend Yields and Expected Stock Returns. \emph{Journal of Financial Economics} 22(1): 3--25.

\hangindent=1.5em\hangafter=1 Fama, E. F. and French, K. R. (1989). Business Conditions and Expected Returns on Stocks and Bonds. \emph{Journal of Financial Economics} 25(1): 23--49.

\hangindent=1.5em\hangafter=1 Fama, E. F. and French, K. R. (1992). The Cross-Section of Expected Stock Returns. \emph{Journal of Finance} 47(2): 427--465.

\hangindent=1.5em\hangafter=1 Fama, E. F. and French, K. R. (1993). Common Risk Factors in the Returns on Stocks and Bonds. \emph{Journal of Financial Economics} 33(1): 3--56.

\hangindent=1.5em\hangafter=1 Fama, E. F. and French, K. R. (1998). Value versus Growth: The International Evidence. \emph{Journal of Finance} 53(6): 1975--1999.

\hangindent=1.5em\hangafter=1 Fama, E. F. and French, K. R. (2008). Dissecting Anomalies. \emph{Journal of Finance} 63(4): 1653--1678.

\hangindent=1.5em\hangafter=1 Fama, E. F. and French, K. R. (2010). Luck versus Skill in the Cross-Section of Mutual Fund Returns. \emph{Journal of Finance} 65(5): 1915--1947.

\hangindent=1.5em\hangafter=1 Fama, E. F. and French, K. R. (2016). Dissecting Anomalies with a Five-Factor Model. \emph{Review of Financial Studies} 29(1): 69--103.

\hangindent=1.5em\hangafter=1 Fama, E. F. and French, K. R. (2021). The Value Premium. \emph{Review of Asset Pricing Studies} 11(1): 105--121.

\hangindent=1.5em\hangafter=1 Fama, E. F. and MacBeth, J. D. (1973). Risk, Return, and Equilibrium: Empirical Tests. \emph{Journal of Political Economy} 81(3): 607--636.

\hangindent=1.5em\hangafter=1 Forbes, K. J. and Rigobon, R. (2002). No Contagion, Only Interdependence: Measuring Stock Market Comovements. \emph{Journal of Finance} 57(5): 2223--2261.

\hangindent=1.5em\hangafter=1 Frazzini, A. and Pedersen, L. H. (2014). Betting against Beta. \emph{Journal of Financial Economics} 111(1): 1--25.

\hangindent=1.5em\hangafter=1 Frazzini, A., Israel, R. and Moskowitz, T. J. (2018). Trading Costs. Working paper, SSRN 3229719.

\hangindent=1.5em\hangafter=1 Freund, Y. and Schapire, R. E. (1997). A Decision-Theoretic Generalization of On-Line Learning and an Application to Boosting. \emph{Journal of Computer and System Sciences} 55(1): 119--139.

\hangindent=1.5em\hangafter=1 Freyberger, J., Neuhierl, A. and Weber, M. (2020). Dissecting Characteristics Nonparametrically. \emph{Review of Financial Studies} 33(5): 2326--2377.

\hangindent=1.5em\hangafter=1 Friedman, J. H. (2001). Greedy Function Approximation: A Gradient Boosting Machine. \emph{Annals of Statistics} 29(5): 1189--1232.

\hangindent=1.5em\hangafter=1 Froot, K. A. and Thaler, R. H. (1990). Anomalies: Foreign Exchange. \emph{Journal of Economic Perspectives} 4(3): 179--192.

\hangindent=1.5em\hangafter=1 Fung, W. and Hsieh, D. A. (2001). The Risk in Hedge Fund Strategies: Theory and Evidence from Trend Followers. \emph{Review of Financial Studies} 14(2): 313--341.

\hangindent=1.5em\hangafter=1 G\^arleanu, N. and Pedersen, L. H. (2013). Dynamic Trading with Predictable Returns and Transaction Costs. \emph{Journal of Finance} 68(6): 2309--2340.

\hangindent=1.5em\hangafter=1 G\^arleanu, N., Pedersen, L. H. and Poteshman, A. M. (2009). Demand-Based Option Pricing. \emph{Review of Financial Studies} 22(10): 4259--4299.

\hangindent=1.5em\hangafter=1 Garman, M. B. and Klass, M. J. (1980). On the Estimation of Security Price Volatilities from Historical Data. \emph{Journal of Business} 53(1): 67--78.

\hangindent=1.5em\hangafter=1 Gelman, A. and Loken, E. (2014). The Statistical Crisis in Science. \emph{American Scientist} 102(6): 460--465.

\hangindent=1.5em\hangafter=1 Geman, S., Bienenstock, E. and Doursat, R. (1992). Neural Networks and the Bias/Variance Dilemma. \emph{Neural Computation} 4(1): 1--58.

\hangindent=1.5em\hangafter=1 Gentzkow, M., Kelly, B. and Taddy, M. (2019). Text as Data. \emph{Journal of Economic Literature} 57(3): 535--574.

\hangindent=1.5em\hangafter=1 Giacomini, R. and White, H. (2006). Tests of Conditional Predictive Ability. \emph{Econometrica} 74(6): 1545--1578.

\hangindent=1.5em\hangafter=1 Gibbons, M. R., Ross, S. A. and Shanken, J. (1989). A Test of the Efficiency of a Given Portfolio. \emph{Econometrica} 57(5): 1121--1152.

\hangindent=1.5em\hangafter=1 Gilchrist, S. and Zakraj\v{s}ek, E. (2012). Credit Spreads and Business Cycle Fluctuations. \emph{American Economic Review} 102(4): 1692--1720.

\hangindent=1.5em\hangafter=1 Glasserman, P. and Lin, C. (2023). Assessing Look-Ahead Bias in Stock Return Predictions Generated by GPT Sentiment Analysis. arXiv:2309.17322.

\hangindent=1.5em\hangafter=1 Glosten, L. R., Jagannathan, R. and Runkle, D. E. (1993). On the Relation between the Expected Value and the Volatility of the Nominal Excess Return on Stocks. \emph{Journal of Finance} 48(5): 1779--1801.

\hangindent=1.5em\hangafter=1 Gneiting, T. (2011). Making and Evaluating Point Forecasts. \emph{Journal of the American Statistical Association} 106(494): 746--762.

\hangindent=1.5em\hangafter=1 Gneiting, T. and Raftery, A. E. (2007). Strictly Proper Scoring Rules, Prediction, and Estimation. \emph{Journal of the American Statistical Association} 102(477): 359--378.

\hangindent=1.5em\hangafter=1 Goel, A., Pasricha, P., Magris, M. and Kanniainen, J. (2025). Foundation Time-Series AI Model for Realized Volatility Forecasting. arXiv:2505.11163.

\hangindent=1.5em\hangafter=1 Goetzmann, W., Ingersoll, J., Spiegel, M. and Welch, I. (2007). Portfolio Performance Manipulation and Manipulation-Proof Performance Measures. \emph{Review of Financial Studies} 20(5): 1503--1546.

\hangindent=1.5em\hangafter=1 Goyal, A. and Welch, I. (2008). A Comprehensive Look at the Empirical Performance of Equity Premium Prediction. \emph{Review of Financial Studies} 21(4): 1455--1508.

\hangindent=1.5em\hangafter=1 Goyal, A., Welch, I. and Zafirov, A. (2024). A Comprehensive 2022 Look at the Empirical Performance of Equity Premium Prediction. \emph{Review of Financial Studies} 37(11): 3490--3557.

\hangindent=1.5em\hangafter=1 Graham, B. and Dodd, D. L. (1934). \emph{Security Analysis}. New York: McGraw-Hill.

\hangindent=1.5em\hangafter=1 Graham, J. R. and Harvey, C. R. (2001). The Theory and Practice of Corporate Finance: Evidence from the Field. \emph{Journal of Financial Economics} 60(2--3): 187--243.

\hangindent=1.5em\hangafter=1 Greenspan, A. (2005). Federal Reserve Board's Semiannual Monetary Policy Report to the Congress. Testimony before the Committee on Banking, Housing, and Urban Affairs, U.S. Senate, February 16.

\hangindent=1.5em\hangafter=1 Grossman, S. J. and Stiglitz, J. E. (1980). On the Impossibility of Informationally Efficient Markets. \emph{American Economic Review} 70(3): 393--408.

\hangindent=1.5em\hangafter=1 Grundy, B. D. and Martin, J. S. (2001). Understanding the Nature of the Risks and the Source of the Rewards to Momentum Investing. \emph{Review of Financial Studies} 14(1): 29--78.

\hangindent=1.5em\hangafter=1 Gu, S., Kelly, B. and Xiu, D. (2020). Empirical Asset Pricing via Machine Learning. \emph{Review of Financial Studies} 33(5): 2223--2273.

\hangindent=1.5em\hangafter=1 Gu, S., Kelly, B. and Xiu, D. (2021). Autoencoder Asset Pricing Models. \emph{Journal of Econometrics} 222(1): 429--450.

\hangindent=1.5em\hangafter=1 G\"urkaynak, R. S., Sack, B. and Wright, J. H. (2007). The U.S. Treasury Yield Curve: 1961 to the Present. \emph{Journal of Monetary Economics} 54(8): 2291--2304.

\hangindent=1.5em\hangafter=1 Hansen, L. P. and Hodrick, R. J. (1980). Forward Exchange Rates as Optimal Predictors of Future Spot Rates: An Econometric Analysis. \emph{Journal of Political Economy} 88(5): 829--853.

\hangindent=1.5em\hangafter=1 Hansen, P. R. (2005). A Test for Superior Predictive Ability. \emph{Journal of Business \& Economic Statistics} 23(4): 365--380.

\hangindent=1.5em\hangafter=1 Hansen, P. R. and Lunde, A. (2005). A Forecast Comparison of Volatility Models: Does Anything Beat a GARCH(1,1)? \emph{Journal of Applied Econometrics} 20(7): 873--889.

\hangindent=1.5em\hangafter=1 Harvey, C. R. (2017). Presidential Address: The Scientific Outlook in Financial Economics. \emph{Journal of Finance} 72(4): 1399--1440.

\hangindent=1.5em\hangafter=1 Harvey, C. R. and Liu, Y. (2015). Backtesting. \emph{Journal of Portfolio Management} 42(1): 13--28.

\hangindent=1.5em\hangafter=1 Harvey, C. R. and Liu, Y. (2020). False (and Missed) Discoveries in Financial Economics. \emph{Journal of Finance} 75(5): 2503--2553.

\hangindent=1.5em\hangafter=1 Harvey, C. R., Liu, Y. and Zhu, H. (2016). \ldots and the Cross-Section of Expected Returns. \emph{Review of Financial Studies} 29(1): 5--68.

\hangindent=1.5em\hangafter=1 Harvey, D., Leybourne, S. and Newbold, P. (1997). Testing the Equality of Prediction Mean Squared Errors. \emph{International Journal of Forecasting} 13(2): 281--291.

\hangindent=1.5em\hangafter=1 Hasbrouck, J. (2009). Trading Costs and Returns for U.S. Equities: Estimating Effective Costs from Daily Data. \emph{Journal of Finance} 64(3): 1445--1477.

\hangindent=1.5em\hangafter=1 Haugen, R. A. and Heins, A. J. (1975). Risk and the Rate of Return on Financial Assets: Some Old Wine in New Bottles. \emph{Journal of Financial and Quantitative Analysis} 10(5): 775--784.

\hangindent=1.5em\hangafter=1 Hodrick, R. J. (1992). Dividend Yields and Expected Stock Returns: Alternative Procedures for Inference and Measurement. \emph{Review of Financial Studies} 5(3): 357--386.

\hangindent=1.5em\hangafter=1 Hoerl, A. E. and Kennard, R. W. (1970). Ridge Regression: Biased Estimation for Nonorthogonal Problems. \emph{Technometrics} 12(1): 55--67.

\hangindent=1.5em\hangafter=1 Holm, S. (1979). A Simple Sequentially Rejective Multiple Test Procedure. \emph{Scandinavian Journal of Statistics} 6(2): 65--70.

\hangindent=1.5em\hangafter=1 Hong, H. and Stein, J. C. (1999). A Unified Theory of Underreaction, Momentum Trading, and Overreaction in Asset Markets. \emph{Journal of Finance} 54(6): 2143--2184.

\hangindent=1.5em\hangafter=1 Hou, K., Xue, C. and Zhang, L. (2020). Replicating Anomalies. \emph{Review of Financial Studies} 33(5): 2019--2133.

\hangindent=1.5em\hangafter=1 Hu, E. J., Shen, Y., Wallis, P., Allen-Zhu, Z., Li, Y., Wang, S., Wang, L. and Chen, W. (2022). LoRA: Low-Rank Adaptation of Large Language Models. \emph{International Conference on Learning Representations}.

\hangindent=1.5em\hangafter=1 Huang, D., Li, J., Wang, L. and Zhou, G. (2020). Time Series Momentum: Is It There? \emph{Journal of Financial Economics} 135(3): 774--794.

\hangindent=1.5em\hangafter=1 Huang, J.-Z. and Huang, M. (2012). How Much of the Corporate-Treasury Yield Spread Is Due to Credit Risk? \emph{Review of Asset Pricing Studies} 2(2): 153--202.

\hangindent=1.5em\hangafter=1 Hull, J. and White, A. (1998). Incorporating Volatility Updating into the Historical Simulation Method for Value-at-Risk. \emph{Journal of Risk} 1(1): 5--19.

\hangindent=1.5em\hangafter=1 Hull, J. and White, A. (2017). Optimal Delta Hedging for Options. \emph{Journal of Banking \& Finance} 82: 180--190.

\hangindent=1.5em\hangafter=1 Hurst, B., Ooi, Y. H. and Pedersen, L. H. (2017). A Century of Evidence on Trend-Following Investing. \emph{Journal of Portfolio Management} 44(1): 15--29.

\hangindent=1.5em\hangafter=1 Inoue, A. and Kilian, L. (2004). In-Sample or Out-of-Sample Tests of Predictability: Which One Should We Use? \emph{Econometric Reviews} 23(4): 371--402.

\hangindent=1.5em\hangafter=1 Israel, R., Kelly, B. and Moskowitz, T. J. (2020). Can Machines ``Learn'' Finance? \emph{Journal of Investment Management} 18(2).

\hangindent=1.5em\hangafter=1 Israelov, R. and Nielsen, L. N. (2014). Covered Call Strategies: One Fact and Eight Myths. \emph{Financial Analysts Journal} 70(6): 23--31.

\hangindent=1.5em\hangafter=1 Israelov, R. and Nielsen, L. N. (2015). Covered Calls Uncovered. \emph{Financial Analysts Journal} 71(6): 44--57.

\hangindent=1.5em\hangafter=1 Jagannathan, R. and Ma, T. (2003). Risk Reduction in Large Portfolios: Why Imposing the Wrong Constraints Helps. \emph{Journal of Finance} 58(4): 1651--1684.

\hangindent=1.5em\hangafter=1 Jegadeesh, N. (1990). Evidence of Predictable Behavior of Security Returns. \emph{Journal of Finance} 45(3): 881--898.

\hangindent=1.5em\hangafter=1 Jegadeesh, N. and Titman, S. (1993). Returns to Buying Winners and Selling Losers: Implications for Stock Market Efficiency. \emph{Journal of Finance} 48(1): 65--91.

\hangindent=1.5em\hangafter=1 Jegadeesh, N. and Titman, S. (2001). Profitability of Momentum Strategies: An Evaluation of Alternative Explanations. \emph{Journal of Finance} 56(2): 699--720.

\hangindent=1.5em\hangafter=1 Jensen, M. C. (1978). Some Anomalous Evidence Regarding Market Efficiency. \emph{Journal of Financial Economics} 6(2--3): 95--101.

\hangindent=1.5em\hangafter=1 Jensen, T. I., Kelly, B. and Pedersen, L. H. (2023). Is There a Replication Crisis in Finance? \emph{Journal of Finance} 78(5): 2465--2518.

\hangindent=1.5em\hangafter=1 Jobson, J. D. and Korkie, B. M. (1981). Performance Hypothesis Testing with the Sharpe and Treynor Measures. \emph{Journal of Finance} 36(4): 889--908.

\hangindent=1.5em\hangafter=1 \emph{Journal of Finance} (2024). Data and Code Sharing Policy, April 1, 2024.

\hangindent=1.5em\hangafter=1 J.P. Morgan/Reuters (1996). \emph{RiskMetrics---Technical Document} (4th ed.). New York: Morgan Guaranty Trust Company.

\hangindent=1.5em\hangafter=1 Kamal, M. and Derman, E. (1999). When You Cannot Hedge Continuously: The Corrections to Black--Scholes. \emph{Risk} 12(1).

\hangindent=1.5em\hangafter=1 Keim, D. B. and Stambaugh, R. F. (1986). Predicting Returns in the Stock and Bond Markets. \emph{Journal of Financial Economics} 17(2): 357--390.

\hangindent=1.5em\hangafter=1 Kelly, B., Kuznetsov, B., Malamud, S. and Xu, T. A. (2025). Artificial Intelligence Asset Pricing Models. NBER Working Paper 33351.

\hangindent=1.5em\hangafter=1 Kelly, B., Kuznetsov, B., Malamud, S., Xu, T. A. and Zhang, Y. (2024). Large and Deep Factor Models. Working paper, arXiv:2402.06635.

\hangindent=1.5em\hangafter=1 Kelly, B., Malamud, S. and Zhou, K. (2024). The Virtue of Complexity in Return Prediction. \emph{Journal of Finance} 79(1): 459--503.

\hangindent=1.5em\hangafter=1 Kelly, B., Pruitt, S. and Su, Y. (2019). Characteristics Are Covariances: A Unified Model of Risk and Return. \emph{Journal of Financial Economics} 134(3): 501--524.

\hangindent=1.5em\hangafter=1 Kendall, M. G. (1954). Note on Bias in the Estimation of Autocorrelation. \emph{Biometrika} 41(3--4): 403--404.

\hangindent=1.5em\hangafter=1 Khandani, A. E. and Lo, A. W. (2011). What Happened to the Quants in August 2007? Evidence from Factors and Transactions Data. \emph{Journal of Financial Markets} 14(1): 1--46.

\hangindent=1.5em\hangafter=1 Khandani, A. E., Kim, A. J. and Lo, A. W. (2010). Consumer Credit-Risk Models via Machine-Learning Algorithms. \emph{Journal of Banking \& Finance} 34(11): 2767--2787.

\hangindent=1.5em\hangafter=1 Kirby, C. and Ostdiek, B. (2012). It's All in the Timing: Simple Active Portfolio Strategies That Outperform Naive Diversification. \emph{Journal of Financial and Quantitative Analysis} 47(2): 437--467.

\hangindent=1.5em\hangafter=1 Korajczyk, R. A. and Sadka, R. (2004). Are Momentum Profits Robust to Trading Costs? \emph{Journal of Finance} 59(3): 1039--1082.

\hangindent=1.5em\hangafter=1 Kosowski, R., Timmermann, A., Wermers, R. and White, H. (2006). Can Mutual Fund ``Stars'' Really Pick Stocks? New Evidence from a Bootstrap Analysis. \emph{Journal of Finance} 61(6): 2551--2595.

\hangindent=1.5em\hangafter=1 Kothari, S. P., Shanken, J. and Sloan, R. G. (1995). Another Look at the Cross-Section of Expected Stock Returns. \emph{Journal of Finance} 50(1): 185--224.

\hangindent=1.5em\hangafter=1 Kozak, S., Nagel, S. and Santosh, S. (2020). Shrinking the Cross-Section. \emph{Journal of Financial Economics} 135(2): 271--292.

\hangindent=1.5em\hangafter=1 Kuester, K., Mittnik, S. and Paolella, M. S. (2006). Value-at-Risk Prediction: A Comparison of Alternative Strategies. \emph{Journal of Financial Econometrics} 4(1): 53--89.

\hangindent=1.5em\hangafter=1 Kuo, C. (2024a). \emph{The Handbook of Anomaly Detection: Cutting-edge Methods and Hands-On Code Examples} (2nd ed.).

\hangindent=1.5em\hangafter=1 Kuo, C. (2024b). \emph{Modern Time Series Forecasting Techniques for Predictive Analytics and Anomaly Detection: From Classical Foundations to Cutting-Edge Applications}.

\hangindent=1.5em\hangafter=1 Kuo, C. (2025). \emph{Introduction to Statistics with Python: The Essential Guide to Data Analysis}.

\hangindent=1.5em\hangafter=1 Kuo, C. (2026). \emph{Agentic AI for Anomaly Detection: Building Your Own Agentic AI Systems Step by Step} (3rd ed.).

\hangindent=1.5em\hangafter=1 Kupiec, P. H. (1995). Techniques for Verifying the Accuracy of Risk Measurement Models. \emph{Journal of Derivatives} 3(2): 73--84.

\hangindent=1.5em\hangafter=1 Kyle, A. S. (1985). Continuous Auctions and Insider Trading. \emph{Econometrica} 53(6): 1315--1335.

\hangindent=1.5em\hangafter=1 Lachanski, M. and Pav, S. (2017). Shy of the Character Limit: ``Twitter Mood Predicts the Stock Market'' Revisited. \emph{Econ Journal Watch} 14(3): 302--345.

\hangindent=1.5em\hangafter=1 Lakonishok, J., Shleifer, A. and Vishny, R. W. (1994). Contrarian Investment, Extrapolation, and Risk. \emph{Journal of Finance} 49(5): 1541--1578.

\hangindent=1.5em\hangafter=1 Ledoit, O. and Wolf, M. (2004). A Well-Conditioned Estimator for Large-Dimensional Covariance Matrices. \emph{Journal of Multivariate Analysis} 88(2): 365--411.

\hangindent=1.5em\hangafter=1 Lehmann, B. N. (1990). Fads, Martingales, and Market Efficiency. \emph{Quarterly Journal of Economics} 105(1): 1--28.

\hangindent=1.5em\hangafter=1 Leinweber, D. J. (2007). Stupid Data Miner Tricks: Overfitting the S\&P 500. \emph{Journal of Investing} 16(1): 15--22.

\hangindent=1.5em\hangafter=1 Leippold, M., Wang, Q. and Zhou, W. (2022). Machine Learning in the Chinese Stock Market. \emph{Journal of Financial Economics} 145(2): 64--82.

\hangindent=1.5em\hangafter=1 Leland, H. E. (1985). Option Pricing and Replication with Transactions Costs. \emph{Journal of Finance} 40(5): 1283--1301.

\hangindent=1.5em\hangafter=1 Lesmond, D. A., Schill, M. J. and Zhou, C. (2004). The Illusory Nature of Momentum Profits. \emph{Journal of Financial Economics} 71(2): 349--380.

\hangindent=1.5em\hangafter=1 Lettau, M. and Van Nieuwerburgh, S. (2008). Reconciling the Return Predictability Evidence. \emph{Review of Financial Studies} 21(4): 1607--1652.

\hangindent=1.5em\hangafter=1 Lewellen, J. (2004). Predicting Returns with Financial Ratios. \emph{Journal of Financial Economics} 74(2): 209--235.

\hangindent=1.5em\hangafter=1 Lintner, J. (1965). The Valuation of Risk Assets and the Selection of Risky Investments in Stock Portfolios and Capital Budgets. \emph{Review of Economics and Statistics} 47(1): 13--37.

\hangindent=1.5em\hangafter=1 Liu, J., Stambaugh, R. F. and Yuan, Y. (2018). Absolving Beta of Volatility's Effects. \emph{Journal of Financial Economics} 128(1): 1--15.

\hangindent=1.5em\hangafter=1 Liu, L. Y., Patton, A. J. and Sheppard, K. (2015). Does Anything Beat 5-Minute RV? A Comparison of Realized Measures across Multiple Asset Classes. \emph{Journal of Econometrics} 187(1): 293--311.

\hangindent=1.5em\hangafter=1 Lo, A. W. (2002). The Statistics of Sharpe Ratios. \emph{Financial Analysts Journal} 58(4): 36--52.

\hangindent=1.5em\hangafter=1 Lo, A. W. and MacKinlay, A. C. (1988). Stock Market Prices Do Not Follow Random Walks: Evidence from a Simple Specification Test. \emph{Review of Financial Studies} 1(1): 41--66.

\hangindent=1.5em\hangafter=1 Lo, A. W. and MacKinlay, A. C. (1990). Data-Snooping Biases in Tests of Financial Asset Pricing Models. \emph{Review of Financial Studies} 3(3): 431--467.

\hangindent=1.5em\hangafter=1 Longin, F. and Solnik, B. (2001). Extreme Correlation of International Equity Markets. \emph{Journal of Finance} 56(2): 649--676.

\hangindent=1.5em\hangafter=1 L\'opez de Prado, M. (2016). Building Diversified Portfolios That Outperform Out of Sample. \emph{Journal of Portfolio Management} 42(4): 59--69.

\hangindent=1.5em\hangafter=1 L\'opez de Prado, M. (2018). \emph{Advances in Financial Machine Learning}. Hoboken, NJ: Wiley.

\hangindent=1.5em\hangafter=1 L\'opez de Prado, M. (2019). A Data Science Solution to the Multiple-Testing Crisis in Financial Research. \emph{Journal of Financial Data Science} 1(1): 99--110.

\hangindent=1.5em\hangafter=1 Lopez-Lira, A. and Tang, Y. (2023). Can ChatGPT Forecast Stock Price Movements? Return Predictability and Large Language Models. arXiv:2304.07619.

\hangindent=1.5em\hangafter=1 Loughran, T. and McDonald, B. (2011). When Is a Liability Not a Liability? Textual Analysis, Dictionaries, and 10-Ks. \emph{Journal of Finance} 66(1): 35--65.

\hangindent=1.5em\hangafter=1 Lundberg, S. M. and Lee, S.-I. (2017). A Unified Approach to Interpreting Model Predictions. \emph{Advances in Neural Information Processing Systems} 30.

\hangindent=1.5em\hangafter=1 Lustig, H., Roussanov, N. and Verdelhan, A. (2011). Common Risk Factors in Currency Markets. \emph{Review of Financial Studies} 24(11): 3731--3777.

\hangindent=1.5em\hangafter=1 Maillard, S., Roncalli, T. and Teiletche, J. (2010). The Properties of Equally Weighted Risk Contribution Portfolios. \emph{Journal of Portfolio Management} 36(4): 60--70.

\hangindent=1.5em\hangafter=1 Mandelbrot, B. (1963). The Variation of Certain Speculative Prices. \emph{Journal of Business} 36(4): 394--419.

\hangindent=1.5em\hangafter=1 Marconi, B. A. (2025). Time Series Foundation Models for Multivariate Financial Time Series Forecasting. arXiv:2507.07296.

\hangindent=1.5em\hangafter=1 Markowitz, H. (1952). Portfolio Selection. \emph{Journal of Finance} 7(1): 77--91.

\hangindent=1.5em\hangafter=1 McLean, R. D. and Pontiff, J. (2016). Does Academic Research Destroy Stock Return Predictability? \emph{Journal of Finance} 71(1): 5--32.

\hangindent=1.5em\hangafter=1 McNeil, A. J. and Frey, R. (2000). Estimation of Tail-Related Risk Measures for Heteroscedastic Financial Time Series: An Extreme Value Approach. \emph{Journal of Empirical Finance} 7(3--4): 271--300.

\hangindent=1.5em\hangafter=1 Meese, R. A. and Rogoff, K. (1983). Empirical Exchange Rate Models of the Seventies: Do They Fit Out of Sample? \emph{Journal of International Economics} 14(1--2): 3--24.

\hangindent=1.5em\hangafter=1 Memmel, C. (2003). Performance Hypothesis Testing with the Sharpe Ratio. \emph{Finance Letters} 1: 21--23.

\hangindent=1.5em\hangafter=1 Menkhoff, L., Sarno, L., Schmeling, M. and Schrimpf, A. (2012). Carry Trades and Global Foreign Exchange Volatility. \emph{Journal of Finance} 67(2): 681--718.

\hangindent=1.5em\hangafter=1 Merton, R. C. (1973). Theory of Rational Option Pricing. \emph{Bell Journal of Economics and Management Science} 4(1): 141--183.

\hangindent=1.5em\hangafter=1 Merton, R. C. (1974). On the Pricing of Corporate Debt: The Risk Structure of Interest Rates. \emph{Journal of Finance} 29(2): 449--470.

\hangindent=1.5em\hangafter=1 Merton, R. C. (1980). On Estimating the Expected Return on the Market: An Exploratory Investigation. \emph{Journal of Financial Economics} 8(4): 323--361.

\hangindent=1.5em\hangafter=1 Michaud, R. O. (1989). The Markowitz Optimization Enigma: Is ``Optimized'' Optimal? \emph{Financial Analysts Journal} 45(1): 31--42.

\hangindent=1.5em\hangafter=1 Miller, E. M. (1977). Risk, Uncertainty, and Divergence of Opinion. \emph{Journal of Finance} 32(4): 1151--1168.

\hangindent=1.5em\hangafter=1 Modigliani, F. and Cohn, R. A. (1979). Inflation, Rational Valuation and the Market. \emph{Financial Analysts Journal} 35(2): 24--44.

\hangindent=1.5em\hangafter=1 Moody's Investors Service (2007). \emph{Corporate Default and Recovery Rates, 1920--2006}. New York: Moody's.

\hangindent=1.5em\hangafter=1 Moreira, A. and Muir, T. (2017). Volatility-Managed Portfolios. \emph{Journal of Finance} 72(4): 1611--1644.

\hangindent=1.5em\hangafter=1 Morningstar (2024). \emph{Mind the Gap 2024}.

\hangindent=1.5em\hangafter=1 Moskowitz, T. J., Ooi, Y. H. and Pedersen, L. H. (2012). Time Series Momentum. \emph{Journal of Financial Economics} 104(2): 228--250.

\hangindent=1.5em\hangafter=1 Mossin, J. (1966). Equilibrium in a Capital Asset Market. \emph{Econometrica} 34(4): 768--783.

\hangindent=1.5em\hangafter=1 Nagel, S. (2012). Evaporating Liquidity. \emph{Review of Financial Studies} 25(7): 2005--2039.

\hangindent=1.5em\hangafter=1 Nagel, S. (2025). Seemingly Virtuous Complexity in Return Prediction. NBER Working Paper 34104.

\hangindent=1.5em\hangafter=1 Nelson, C. R. and Kim, M. J. (1993). Predictable Stock Returns: The Role of Small Sample Bias. \emph{Journal of Finance} 48(2): 641--661.

\hangindent=1.5em\hangafter=1 Newey, W. K. and West, K. D. (1987). A Simple, Positive Semi-Definite, Heteroskedasticity and Autocorrelation Consistent Covariance Matrix. \emph{Econometrica} 55(3): 703--708.

\hangindent=1.5em\hangafter=1 Nosek, B. A., Ebersole, C. R., DeHaven, A. C. and Mellor, D. T. (2018). The Preregistration Revolution. \emph{Proceedings of the National Academy of Sciences} 115(11): 2600--2606.

\hangindent=1.5em\hangafter=1 Novy-Marx, R. (2013). The Other Side of Value: The Gross Profitability Premium. \emph{Journal of Financial Economics} 108(1): 1--28.

\hangindent=1.5em\hangafter=1 Novy-Marx, R. and Velikov, M. (2016). A Taxonomy of Anomalies and Their Trading Costs. \emph{Review of Financial Studies} 29(1): 104--147.

\hangindent=1.5em\hangafter=1 Novy-Marx, R. and Velikov, M. (2022). Betting against Betting against Beta. \emph{Journal of Financial Economics} 143(1): 80--106.

\hangindent=1.5em\hangafter=1 Odean, T. (1998). Are Investors Reluctant to Realize Their Losses? \emph{Journal of Finance} 53(5): 1775--1798.

\hangindent=1.5em\hangafter=1 Page, S. and Panariello, R. A. (2018). When Diversification Fails. \emph{Financial Analysts Journal} 74(3): 19--32.

\hangindent=1.5em\hangafter=1 Pardo, R. E. (1992). \emph{Design, Testing, and Optimization of Trading Systems}. New York: John Wiley \& Sons.

\hangindent=1.5em\hangafter=1 Pardo, R. (2008). \emph{The Evaluation and Optimization of Trading Strategies} (2nd ed.). Hoboken, NJ: John Wiley \& Sons.

\hangindent=1.5em\hangafter=1 Parkinson, M. (1980). The Extreme Value Method for Estimating the Variance of the Rate of Return. \emph{Journal of Business} 53(1): 61--65.

\hangindent=1.5em\hangafter=1 Patton, A. J. (2011). Volatility Forecast Comparison Using Imperfect Volatility Proxies. \emph{Journal of Econometrics} 160(1): 246--256.

\hangindent=1.5em\hangafter=1 P\'erignon, C. and Smith, D. R. (2010). The Level and Quality of Value-at-Risk Disclosure by Commercial Banks. \emph{Journal of Banking \& Finance} 34(2): 362--377.

\hangindent=1.5em\hangafter=1 Perold, A. F. (1988). The Implementation Shortfall: Paper versus Reality. \emph{Journal of Portfolio Management} 14(3): 4--9.

\hangindent=1.5em\hangafter=1 Petersen, M. A. (2009). Estimating Standard Errors in Finance Panel Data Sets: Comparing Approaches. \emph{Review of Financial Studies} 22(1): 435--480.

\hangindent=1.5em\hangafter=1 Piotroski, J. D. (2000). Value Investing: The Use of Historical Financial Statement Information to Separate Winners from Losers. \emph{Journal of Accounting Research} 38(Supplement): 1--41.

\hangindent=1.5em\hangafter=1 Politis, D. N. and Romano, J. P. (1994). The Stationary Bootstrap. \emph{Journal of the American Statistical Association} 89(428): 1303--1313.

\hangindent=1.5em\hangafter=1 Pontiff, J. and Woodgate, A. (2008). Share Issuance and Cross-Sectional Returns. \emph{Journal of Finance} 63(2): 921--945.

\hangindent=1.5em\hangafter=1 Poon, S.-H. and Granger, C. W. J. (2003). Forecasting Volatility in Financial Markets: A Review. \emph{Journal of Economic Literature} 41(2): 478--539.

\hangindent=1.5em\hangafter=1 Presidential Task Force on Market Mechanisms (1988). \emph{Report of the Presidential Task Force on Market Mechanisms} (the Brady Report). Washington, DC: U.S. Government Printing Office.

\hangindent=1.5em\hangafter=1 Rahimikia, E., Ni, H. and Wang, W. (2025). Re(visiting) Time Series Foundation Models in Finance. arXiv:2511.18578.

\hangindent=1.5em\hangafter=1 Rapach, D. E., Strauss, J. K. and Zhou, G. (2010). Out-of-Sample Equity Premium Prediction: Combination Forecasts and Links to the Real Economy. \emph{Review of Financial Studies} 23(2): 821--862.

\hangindent=1.5em\hangafter=1 Roll, R. (1984). A Simple Implicit Measure of the Effective Bid-Ask Spread in an Efficient Market. \emph{Journal of Finance} 39(4): 1127--1139.

\hangindent=1.5em\hangafter=1 Rosenberg, B., Reid, K. and Lanstein, R. (1985). Persuasive Evidence of Market Inefficiency. \emph{Journal of Portfolio Management} 11(3): 9--16.

\hangindent=1.5em\hangafter=1 Rouwenhorst, K. G. (1998). International Momentum Strategies. \emph{Journal of Finance} 53(1): 267--284.

\hangindent=1.5em\hangafter=1 Rubinstein, M. (1994). Implied Binomial Trees. \emph{Journal of Finance} 49(3): 771--818.

\hangindent=1.5em\hangafter=1 Samuelson, P. A. (1965). Proof That Properly Anticipated Prices Fluctuate Randomly. \emph{Industrial Management Review} 6(2): 41--49.

\hangindent=1.5em\hangafter=1 Shanken, J. (1992). On the Estimation of Beta-Pricing Models. \emph{Review of Financial Studies} 5(1): 1--33.

\hangindent=1.5em\hangafter=1 Sharpe, W. F. (1964). Capital Asset Prices: A Theory of Market Equilibrium under Conditions of Risk. \emph{Journal of Finance} 19(3): 425--442.

\hangindent=1.5em\hangafter=1 Sharpe, W. F. (1966). Mutual Fund Performance. \emph{Journal of Business} 39(1): 119--138.

\hangindent=1.5em\hangafter=1 Shefrin, H. and Statman, M. (1985). The Disposition to Sell Winners Too Early and Ride Losers Too Long: Theory and Evidence. \emph{Journal of Finance} 40(3): 777--790.

\hangindent=1.5em\hangafter=1 Shi, Y., Fu, Z., Chen, S., Zhao, B., Xu, W., Zhang, C. and Li, J. (2025). Kronos: A Foundation Model for the Language of Financial Markets. arXiv:2508.02739.

\hangindent=1.5em\hangafter=1 Shiller, R. J. (2000). \emph{Irrational Exuberance}. Princeton, NJ: Princeton University Press.

\hangindent=1.5em\hangafter=1 Shumway, T. (1997). The Delisting Bias in CRSP Data. \emph{Journal of Finance} 52(1): 327--340.

\hangindent=1.5em\hangafter=1 Shumway, T. and Warther, V. A. (1999). The Delisting Bias in CRSP's Nasdaq Data and Its Implications for the Size Effect. \emph{Journal of Finance} 54(6): 2361--2379.

\hangindent=1.5em\hangafter=1 Simonsohn, U., Simmons, J. P. and Nelson, L. D. (2020). Specification Curve Analysis. \emph{Nature Human Behaviour} 4(11): 1208--1214.

\hangindent=1.5em\hangafter=1 Sloan, R. G. (1996). Do Stock Prices Fully Reflect Information in Accruals and Cash Flows about Future Earnings? \emph{The Accounting Review} 71(3): 289--315.

\hangindent=1.5em\hangafter=1 Stambaugh, R. F. (1999). Predictive Regressions. \emph{Journal of Financial Economics} 54(3): 375--421.

\hangindent=1.5em\hangafter=1 Stambaugh, R. F., Yu, J. and Yuan, Y. (2012). The Short of It: Investor Sentiment and Anomalies. \emph{Journal of Financial Economics} 104(2): 288--302.

\hangindent=1.5em\hangafter=1 Sullivan, R., Timmermann, A. and White, H. (1999). Data-Snooping, Technical Trading Rule Performance, and the Bootstrap. \emph{Journal of Finance} 54(5): 1647--1691.

\hangindent=1.5em\hangafter=1 Swiss Federal Banking Commission (2008). \emph{Subprime Crisis: SFBC Investigation into the Causes of the Write-downs of UBS AG}. Bern: SFBC.

\hangindent=1.5em\hangafter=1 Tashman, L. J. (2000). Out-of-Sample Tests of Forecasting Accuracy: An Analysis and Review. \emph{International Journal of Forecasting} 16(4): 437--450.

\hangindent=1.5em\hangafter=1 Tech Brew (2022). How the World's First AI-Managed ETF Stacks Up Almost Five Years Later. January 26.

\hangindent=1.5em\hangafter=1 Tetlock, P. C. (2007). Giving Content to Investor Sentiment: The Role of Media in the Stock Market. \emph{Journal of Finance} 62(3): 1139--1168.

\hangindent=1.5em\hangafter=1 Thornton, D. L. and Valente, G. (2012). Out-of-Sample Predictions of Bond Excess Returns and Forward Rates: An Asset Allocation Perspective. \emph{Review of Financial Studies} 25(10): 3141--3168.

\hangindent=1.5em\hangafter=1 Tibshirani, R. (1996). Regression Shrinkage and Selection via the Lasso. \emph{Journal of the Royal Statistical Society, Series B} 58(1): 267--288.

\hangindent=1.5em\hangafter=1 T\'oth, B., Lemp\'eri\`ere, Y., Deremble, C., de Lataillade, J., Kockelkoren, J. and Bouchaud, J.-P. (2011). Anomalous Price Impact and the Critical Nature of Liquidity in Financial Markets. \emph{Physical Review X} 1(2): 021006.

\hangindent=1.5em\hangafter=1 UBS AG (2008). \emph{Shareholder Report on UBS's Write-Downs}. Zurich: UBS.

\hangindent=1.5em\hangafter=1 U.S. Securities and Exchange Commission (2021). SEC Charges LJM Funds Management and Portfolio Managers (complaint), June 2021.

\hangindent=1.5em\hangafter=1 Valkanov, R. (2003). Long-Horizon Regressions: Theoretical Results and Applications. \emph{Journal of Financial Economics} 68(2): 201--232.

\hangindent=1.5em\hangafter=1 Welch, I. (ed.). \emph{Critical Finance Review}, replication issues. \url{https://cfr.ivo-welch.info}.

\hangindent=1.5em\hangafter=1 West, K. D. (1996). Asymptotic Inference about Predictive Ability. \emph{Econometrica} 64(5): 1067--1084.

\hangindent=1.5em\hangafter=1 Whaley, R. E. (2002). Return and Risk of CBOE Buy Write Monthly Index. \emph{Journal of Derivatives} 10(2): 35--42.

\hangindent=1.5em\hangafter=1 White, H. (1980). A Heteroskedasticity-Consistent Covariance Matrix Estimator and a Direct Test for Heteroskedasticity. \emph{Econometrica} 48(4): 817--838.

\hangindent=1.5em\hangafter=1 White, H. (2000). A Reality Check for Data Snooping. \emph{Econometrica} 68(5): 1097--1126.

\hangindent=1.5em\hangafter=1 Zhang, L. (2005). The Value Premium. \emph{Journal of Finance} 60(1): 67--103.

\hangindent=1.5em\hangafter=1 Zillow Group (2021). Form 8-K, November 2, 2021.

\hangindent=1.5em\hangafter=1 Zou, H. and Hastie, T. (2005). Regularization and Variable Selection via the Elastic Net. \emph{Journal of the Royal Statistical Society, Series B} 67(2): 301--320.

\endgroup
